\documentclass[11pt]{article}
\usepackage[T1]{fontenc}
\usepackage[margin=1in]{geometry}
\usepackage{amsmath,amssymb,amsthm,booktabs,array,longtable}
\usepackage[colorlinks=true,linkcolor=blue,citecolor=blue,urlcolor=blue]{hyperref}
\usepackage{microtype}
\newtheorem{theorem}{Theorem}
\newtheorem{corollary}[theorem]{Corollary}
\newcommand{\fl}[1]{\left\lfloor #1\right\rfloor}
\newcommand{\leanurl}{https://github.com/the-omega-institute/trureturing/blob/b95cdbb8cccca425964251ba38cc0b0e517e478a/D5/S3/Arith/FactorialRatio/BalaA295462Parity.lean}
\title{Parity and product extensions of Bala's divisibilities (2026): OEIS A295462, References (record \#19; contribution dated 23 August 2026)}
\author{Haobo Ma\\\small ChronoAI Pte Ltd, Singapore
\and Wenlin Zhang\\\small National University of Singapore\\\small ChronoAI Pte Ltd, Singapore}
\date{4 November 2026}
\begin{document}
\maketitle
\begin{abstract}
Extending Bala's divisibility results (Bala, 2026), we prove his
conjecture in OEIS A295462, References (\%D),
record \#19; contribution dated 23 August 2026:\newline the integer
$a(n)=(15n)!(2n)!/((10n)!(6n)!n!)$ is odd exactly when consecutive
one-bits of $n$ have at least three zero-bits between them.
A nonnegative floor function expresses the prime valuations of $a(n)$;
the lowest pair of one-bits at distance at most three forces a positive
binary valuation summand. We reprove Bala's two divisibilities
(Bala, 2026) and deduce the same parity criterion for
$a(n)/((6n+1)(10n+1))$. The binary description gives written alternative
proofs of the eleven XOR identities in A048718, verified by
Karlsson (2022) with Walnut. For the related ratio
in A295431 we prove the five fixed divisibilities and give explicit
constants for both product-divisibility families posed on 28 August 2025.
\end{abstract}

\section{The conjecture and its setting}
Throughout, $n$ is a nonnegative integer. Define
\begin{equation}\label{eq:a}
 a(n)=\frac{(15n)!(2n)!}{(10n)!(6n)!n!}.
\end{equation}
The References field of A295462, internal record \#19, attributes the
binary-spacing conjecture to Peter Bala on 23 August 2026 \cite{oeis}.
The sequence itself was entered by Gheorghe Coserea in 2017.
Write $n=\sum_{i\in S(n)}2^i$, where $S(n)$ is its finite binary support.
The support is \emph{four-separated} if $i<j$ in $S(n)$ implies $j-i\geq4$;
the empty support is four-separated. These are the integers in A048718
\cite{spaced}.

\begin{theorem}\label{thm:main}
For every $n\geq0$, $a(n)$ is a positive integer, and
\[
 a(n)\equiv1\pmod2
 \quad\Longleftrightarrow\quad S(n)\text{ is four-separated}.
\]
\end{theorem}

The factorial ratio is the sporadic parameter pair $[15,2];[10,6,1]$
in Bober's Table~2, row~32 \cite{bober}. The sporadic-parameter list
linked from A295431 includes this pair \cite{sporadic}; A295431 itself
is the different ratio considered in Section~5. Integrality of both ratios precedes
the conjecture. Soundararajan later gave an elementary proof of the
classification \cite{sound}. Bober's Theorem~3.1 and Lemma~3.2 explain the general
floor-function criterion. Bala's three-page note of 23 August 2026
proves divisibility by $6n+1$ and $10n+1$ \cite{bala}.
The proof below supplies the binary characterization and treats those
divisibilities with the same floor function.

\section{The floor function and the binary proof}
Put
\begin{equation}\label{eq:F}
 F(x)=\fl{15x}+\fl{2x}-\fl{10x}-\fl{6x}-\fl{x}.
\end{equation}
Since $15+2=10+6+1$, $F(x+1)=F(x)$. All its discontinuities
on $[0,1)$ lie at multiples of $1/30$. On
$j/30\leq x<(j+1)/30$, its value is
$\fl{j/2}+\fl{j/15}-\fl{j/3}-\fl{j/5}$:
\begin{center}
\begin{tabular}{c|rrrrrrrrrr}
\toprule
$j$&0&1&2&3&4&5&6&7&8&9\\
$F$&0&0&1&0&1&0&0&0&1&0\\
\midrule
$j$&10&11&12&13&14&15&16&17&18&19\\
$F$&0&0&0&0&1&0&1&1&1&1\\
\midrule
$j$&20&21&22&23&24&25&26&27&28&29\\
$F$&1&0&1&1&1&0&1&0&1&1\\
\bottomrule
\end{tabular}
\end{center}
In particular, $0\leq F(x)\leq1$ for $x\geq0$, including all
endpoints by the half-open interval convention.
Counting multiples of $p,p^2,\ldots$ among the factors of a factorial gives
Legendre's formula. Hence for every prime $p$,
\begin{equation}\label{eq:valuation}
 v_p(a(n))=\sum_{k\geq1}F(n/p^k)\geq0.
\end{equation}
This is a finite sum, because every term with $p^k>15n$ vanishes.
The positive rational number \eqref{eq:a} therefore has no prime in its
reduced denominator and is an integer. For $n=0$ it equals~1.
It is odd if and only if every summand for $p=2$ is zero.

The binary estimate needed in both directions is simple. If the occupied
positions below an occupied position $e$ are four-separated and at distance
at least four from $e$, their contribution $R$ satisfies
\begin{equation}\label{eq:tail}
 0\leq R\leq\sum_{j\geq1,\ e-4j\geq0}2^{e-4j}
       <2^e\sum_{j\geq1}16^{-j}=\frac{2^e}{15}.
\end{equation}
The bound follows by ordering the positions downward: the $j$th is
at most $e-4j$. Strictness follows from finiteness, including $R=0$.

Suppose first that $S(n)$ is four-separated. Fix $k\geq1$ and set
$r=n\bmod2^k$. Periodicity gives $F(n/2^k)=F(r/2^k)$.
If $r=0$ the latter is zero. Otherwise write
$r=2^t+R$, with $t<k$ its highest occupied position.
Then $0\leq15R<2^t$ by \eqref{eq:tail}. Write $\ell=k-t$.
The following table lists all five floors in \eqref{eq:F} at
$x=(2^t+R)/2^k$:
\begin{center}
\begin{tabular}{c|rrrrr|r}
\toprule
$\ell$&$\fl{15x}$&$\fl{2x}$&$\fl{10x}$&$\fl{6x}$&$\fl{x}$&$F(x)$\\
\midrule
1&7&1&5&3&0&0\\
2&3&0&2&1&0&0\\
3&1&0&1&0&0&0\\
$\geq4$&0&0&0&0&0&0\\
\bottomrule
\end{tabular}
\end{center}
Indeed $x=(1+u)/2^\ell$ with $0\leq u<1/15$;
substitution gives the first three rows, while $15x<1$ for $\ell\geq4$.
Thus every binary summand in \eqref{eq:valuation} vanishes.

Conversely, suppose the support is not four-separated. Choose the smallest
occupied position $t$ having an occupied predecessor within distance three,
and let $e$ be the occupied position immediately below $t$.
Then $d=t-e\in\{1,2,3\}$, and all occupied positions below $t$
are four-separated by the minimality of $t$.
Consequently
\[
 n\bmod2^{t+1}=2^t+2^e+R=(2^d+1)2^e+R,
 \qquad 0\leq15R<2^e.
\]
At $x=(n\bmod2^{t+1})/2^{t+1}$ the floors are as follows:
\begin{center}
\begin{tabular}{c|rrrrr|r}
\toprule
$d$&$\fl{15x}$&$\fl{2x}$&$\fl{10x}$&$\fl{6x}$&$\fl{x}$&$F(x)$\\
\midrule
1&11&1&7&4&0&1\\
2&9&1&6&3&0&1\\
3&8&1&5&3&0&1\\
\bottomrule
\end{tabular}
\end{center}
Here $x=(2^d+1+u)/2^{d+1}$ with $0\leq u<1/15$, so each
row follows by direct substitution. Therefore
$F(n/2^{t+1})=1$. Every other summand is nonnegative;
\eqref{eq:valuation} gives $v_2(a(n))\geq1$.
This proves Theorem~\ref{thm:main}.

\section{Divisibility and consequences}
For completeness we give Bala's divisibility argument \cite{bala}, including
the four cases for $10n+1$.
\begin{corollary}\label{cor:quotient}
The number
\[
 b(n)=\frac{a(n)}{(6n+1)(10n+1)}
\]
is a positive integer for all $n\geq0$. It is odd exactly when $S(n)$
is four-separated.
\end{corollary}
\begin{proof}
Let $h$ be 6 or 10. If $p\mid hn+1$ is prime, then $p$ is coprime
to $h$. For $1\leq r\leq v_p(hn+1)$, put $m=p^r$ and write
$s=n\bmod m$. We have $hs=qm-1$ with $1\leq q<h$ and
$\gcd(q,h)=1$. Thus $s/m=q/h-1/(hm)$.
For $h=6$ we have $m\geq5$ and $q=1,5$; for $h=10$ we have
$m\geq3$ and $q=1,3,7,9$. Substitution in \eqref{eq:F} gives
\begin{center}
\begin{tabular}{cc|rrrrr|r}
\toprule
$h$&$q$&$\fl{15s/m}$&$\fl{2s/m}$&$\fl{10s/m}$&$\fl{6s/m}$&$\fl{s/m}$&$F(s/m)$\\
\midrule
6&1&2&0&1&0&0&1\\
6&5&12&1&8&4&0&1\\
10&1&1&0&0&0&0&1\\
10&3&4&0&2&1&0&1\\
10&7&10&1&6&4&0&1\\
10&9&13&1&8&5&0&1\\
\bottomrule
\end{tabular}
\end{center}
For example, the first column is $\fl{15q/h-15/(hm)}$;
the stated lower bounds on $m$ give each row, including equality at
the smallest endpoint. Periodicity and \eqref{eq:valuation} now imply
$v_p(a(n))\geq v_p(hn+1)$. Hence both $6n+1$ and $10n+1$ divide
$a(n)$. Their greatest common divisor divides
$5(6n+1)-3(10n+1)=2$ and is odd, so it is~1.
Their product divides $a(n)$ and is odd. Consequently $v_2(b(n))=v_2(a(n))$,
which proves the parity assertion.
\end{proof}

\begin{corollary}\label{cor:count}
Let $c_m$ count the odd values $a(n)$ with $0\leq n<2^m$. Then
\[
 c_0=1,\quad c_1=2,\quad c_2=3,\quad c_3=4,\qquad
 c_m=c_{m-1}+c_{m-4}\quad(m\geq4).
\]
Let $u_0,u_1,\ldots$ be their indices in increasing order. For $j>0$, let
$e$ be the unique integer with $c_e\leq j<c_{e+1}$. Then
$u_j=2^e+u_{j-c_e}$, with $u_0=0$.
\end{corollary}
\begin{proof}
Represent every $n$ by a word of length $m$,
allowing leading zeros. A four-separated word either starts with 0,
leaving an arbitrary such word of length $m-1$, or starts with 1000,
leaving one of length $m-4$. The initial counts include the all-zero word.
There are exactly $c_e$ admissible indices below $2^e$. Those in
$[2^e,2^{e+1})$ are $2^e$ plus the admissible suffixes below $2^{e-3}$
when $e\geq3$; for $e<3$ the suffix is zero. This proves the ordered recursion.
\end{proof}

\section{The eleven XOR identities in A048718}
Write $\oplus$, $\mathbin{\&}$ and $\mathbin{|}$ for bitwise XOR,
AND and OR. Hanna listed the following identities on 22 January 2006;
Karlsson supplied Walnut formulas on 31 December 2022
\cite{spaced,walnut}. We give written proofs from the binary support.

\begin{corollary}\label{cor:xor}
If $a(n)$ is odd, then
\[
 (un)\oplus(vn)=(u\oplus v)n\quad(0\leq u,v<16).
\]
The same formula holds with XOR replaced throughout by AND or OR.
In particular, $(un)\oplus(vn)=wn$ for each of the eleven triples
\[
\begin{gathered}
(1,14,15),\ (3,9,10),\ (3,13,14),\ (5,9,12),\ (5,11,14),\\
(6,11,13),\ (7,9,14),\ (7,10,13),\ (7,11,12),\\
(12,21,25),\ (12,37,41).
\end{gathered}
\]
\end{corollary}
\begin{proof}
Theorem~\ref{thm:main} makes $S=S(n)$ four-separated. For $u<16$,
the sum $un=\sum_{e\in S}u2^e$ occupies disjoint four-bit blocks.
The three bitwise operations therefore act separately in each block.
This proves the first nine triples and the two analogous formulas.

For the last two triples take $k=4$ or $5$ and put $z=n+2^kn$.
The supports $S\cup(S+k)$ avoid $(S+2)\cup(S+3)$. A collision in
the addition defining $z$ occurs at $t\in S\cap(S+k)$ and creates
a carry at $t+1$. That position belongs to neither $S$ nor $S+k$:
otherwise it would be adjacent to $t$ or to $t-k$ in $S$.
Thus all carries stop after one position. The carry positions also
avoid $S+2$ and $S+3$, since $t\in S$. Consequently $z$, $4n$ and
$8n$ have pairwise disjoint supports, and
\[
 (z+4n)\oplus(z+8n)=(4n)\oplus(8n)=12n.
\]
The pairs $(z+4n,z+8n)$ are $(21n,25n)$ and $(37n,41n)$.
XORing once more with $12n$ gives the displayed triples.
The same disjoint-block argument works for $r$-separated support and
masks $u,v<2^r$, for every $r\geq1$.
\end{proof}

\section{The divisibility questions in A295431}
The directly linked entry A295431 concerns
\[
 T(n)=\frac{(12n)!n!}{(6n)!(4n)!(3n)!}.
\]
Its Comments field (record \#101) contains Bala's contribution dated
28 August 2025, which poses four
divisibilities involving $n+1,2n+1,3n+1$ and a fifth involving
$12n-1$, followed by two product families \cite{sporadic}.
For $R\geq1$, let $L_R=\operatorname{lcm}(1,2,\ldots,R)$.

\begin{corollary}\label{cor:families}
For every $n\geq0$, each of
\[
 \frac{385T(n)}{n+1},\quad \frac{5T(n)}{2n+1},\quad
 \frac{T(n)}{3n+1},\quad
 \frac{770T(n)}{(n+1)(2n+1)(3n+1)},\quad
 \frac{T(n)}{12n-1}
\]
is an integer. More generally the choices
\[
 C(k,R)=36(R-1)!L_{12R}\quad(k=1,2,3),\qquad
 D(R)=(R-1)!L_R
\]
make both expressions
\[
 \frac{C(k,R)T(n)}{\prod_{i=1}^R(kn+i)},\qquad
 \frac{D(R)T(n)}{\prod_{1\leq i\leq R,\ \gcd(i,12)=1}(12n-i)}
\]
integers for every $n\geq0$ and $R\geq1$.
\end{corollary}
\begin{proof}
Define $G(x)=\fl{12x}+\fl{x}-\fl{6x}-\fl{4x}-\fl{3x}$.
It has period one; on the twelve successive half-open cells of
length $1/12$ its values are
\[
 0,1,1,1,0,1,0,1,0,0,0,1.
\]
Thus $T(n)$ is an integer and $v_p(T(n))=\sum_{j\geq1}G(n/p^j)$.
If $s_p(m)$ denotes the sum of the base-$p$ digits of $m$,
Legendre's formula also gives
\begin{equation}\label{eq:Tdigits}
 v_2(T(n))=s_2(3n),\qquad v_3(T(n))=s_3(2n)/2.
\end{equation}

First consider the five fixed assertions. At a prime-power level
$m=p^j$ dividing $hn+1$, reduce $n$ modulo $m$ to obtain
$n/m\equiv q/h-1/(hm)\pmod1$. For $p\geq5$, the cells above give
\[
\begin{array}{c|c|c}
h&\text{allowed }q&G=0\text{ levels}\\\hline
1&1&m=5,7,11\\
2&1&m=5\\
3&1,2&\text{none}
\end{array}
\]
At all other allowed levels $G=1$. For example, the three expressions
approach $1,1/2$ or $1/3,2/3$ from the left; substitution checks
the finitely many levels below 12. The powers $m\geq12$ remain
in the adjacent cell with value one.
For $p=2$, the binary tails of $3n$ give
$s_2(3n)\geq v_2(n+1),v_2(3n+1)$:
if $n+1=2^ra$ with $r\geq2$, the last $r$ digits of $3n$
are those of $2^r-3$, of weight $r-1$, and the positive prefix
$3a-1$ supplies another one; $r=1$ is immediate.
If $2^r\mid3n+1$, the last $r$ bits of $3n$ are all one.
For $p=3$, if $n+1=3^ra$, the last $r$ ternary digits of $2n$
have sum $2r-1$ and the positive prefix $2a-1$ adds at least one.
If $3^r\mid2n+1$, those digits are all two. Equation~\eqref{eq:Tdigits}
therefore supplies both required ternary bounds.
The first three assertions follow. The three denominators are
pairwise coprime except that $\gcd(n+1,3n+1)$ divides 2.
When both are even one of their binary valuations is one, so their
total is at most $v_2(T(n))+1$. The missing $5$-levels for
$n+1$ and $2n+1$ occur in distinct residue classes modulo 5.
This proves the assertion with 770.
For $m\mid12n-1$, necessarily $p\geq5$ and
$n/m\equiv q/12+1/(12m)$, with $q=1,5,7,11$.
The least allowed $m$ in those four classes is respectively
$11,7,5,13$, and each expression lies in a cell with $G=1$.
This proves the fifth assertion, including its negative value at $n=0$.

For the general families we use the following elementary bound.
For any nonzero integers $kn+i$ with $1\leq i\leq R$, or any subset
of the integers $12n-i$ in that range, put $e$ equal to the largest
of their $p$-adic valuations. Then
\begin{equation}\label{eq:productbound}
 v_p(\text{product})\leq e+v_p((R-1)!).
\end{equation}
Indeed, choose an index $i_0$ of largest valuation. For $i\ne i_0$,
the valuation of that factor is at most $v_p(i-i_0)$, by subtracting
the two factors. Their product of differences divides
$(i_0-1)!(R-i_0)!$, which divides $(R-1)!$.

Put $E_p=\fl{\log_p(12R)}$. For $p\geq5$ and
$p^j\mid kn+i$ with $p^j>12R$, reduction modulo $p^j$ gives
\[
 \frac{n\bmod p^j}{p^j}=\frac{q}{k}-\frac{i}{kp^j},
 \qquad 1\leq q\leq k.
\]
The left-hand neighborhoods of width $1/12$ at
$1/3,1/2,2/3,1$ all have $G=1$. Thus $v_p(T(n))\geq e-E_p$.
For $p=2,3$ the corresponding bound is
\begin{equation}\label{eq:smallprimes}
 v_p(T(n))\geq e-E_p-2.
\end{equation}
Here are the digit details. If $z\equiv-d\pmod{p^j}$ with
$0<d<p^j$, then
\[
 s_p(z)\geq s_p(p^j-d)=(p-1)j-s_p(d-1).
\]
For $p=2$, the congruences for $k=1,3$ give respectively
$3n\equiv-3i,-i\pmod{2^j}$. For $p=3$, those for $k=1,2$
give $2n\equiv-2i,-i\pmod{3^j}$. In each case $d\leq3R$;
if $p^j\leq d$ then $j\leq E_p$, and otherwise
$s_p(d-1)\leq(p-1)(E_p+1)$, giving the bound with $E_p+1$.
For the remaining case $k=p$, either $p\nmid i$ and the valuation
is zero, or write $i=ph$ and use
$v_p(pn+i)=1+v_p(n+h)$, losing at most one more.
This proves \eqref{eq:smallprimes}.
Combining these bounds with \eqref{eq:productbound} proves the
first family, since $v_p(L_{12R})=E_p$ and $36=2^2 3^2$.

In the second family every denominator factor is coprime to 6.
If $m=p^j>R$ divides $12n-i$, then
\[
 \frac{n\bmod m}{m}=\frac{q}{12}+\frac{i}{12m},
 \qquad q\in\{1,5,7,11\}.
\]
Each right-hand neighborhood of width $1/12$ at these four
fractions has $G=1$. Therefore $v_p(T(n))\geq
e-\fl{\log_p R}$. Equation~\eqref{eq:productbound} and
$v_p(L_R)=\fl{\log_p R}$ prove the claimed choice of $D(R)$.
No factor $12n-i$ is zero because $i$ is coprime to 12.
\end{proof}

\section{Other source formulas and coverage}
\begin{corollary}\label{cor:formulas}
The binomial identities, hypergeometric series and asymptotic formulas
in A295462 and A295431 hold, as does Mathar's recurrence for A295431.
\end{corollary}
\begin{proof}
The two binomial identities credited to Chai Wah Wu in A295462 follow
directly by expanding the factorials:
\[
 a(n)=\frac{\binom{15n}{5n}\binom{2n}{n}}{\binom{6n}{n}}
     =\frac{\binom{15n}{5n}\binom{7n}{n}}{\binom{7n}{2n}}.
\]
The entry's generating function can likewise be read from the coefficient
ratio. For
\[
 A=\{1,2,4,7,8,11,13,14\}/15,\quad
 B=\{1/10,1/6,3/10,1/2,7/10,5/6,9/10\},\quad
 C=\frac{61509375}{16384},
\]
factorial cancellation gives
\[
 \frac{a(n+1)}{a(n)}
 =\frac{C\prod_{\alpha\in A}(n+\alpha)}
 {(n+1)\prod_{\beta\in B}(n+\beta)}.
\]
The initial coefficient is $a(0)=1$, so the defining coefficient formula
for a hypergeometric series yields
$\sum_{n\geq0}a(n)z^n={}_8F_7(A;B;Cz)$ as a formal power series.
The value of $C$ follows from
$15^{15}2^2/(10^{10}6^6)=61509375/16384$.
Finally, Stirling's formula applied to the five factorials gives
\[
 a(n)\sim\frac{1}{2\sqrt{\pi n}}
 \left(\frac{3^9 5^5}{2^{14}}\right)^n
 \qquad(n\longrightarrow\infty),
\]
as recorded by Amiram Eldar in A295462. Indeed the exponential factors
$e^{-kn}$ cancel, the powers of $n$ leave $n^{-1/2}$, and the square-root
factor is $\sqrt{30/60}/\sqrt{2\pi n}=1/(2\sqrt{\pi n})$.

For $T(n)$ the same cancellations give Wu's two formulas in A295431:
\[
 T(n)=\frac{\binom{12n}{6n}\binom{6n}{2n}}{\binom{3n}{n}}
     =\frac{\binom{12n}{6n}\binom{7n}{3n}}{\binom{7n}{n}}.
\]
Its consecutive ratio, after removing the common linear factors, is
\[
 \frac{T(n)}{T(n-1)}=
 \frac{24(12n-11)(12n-1)(12n-5)(12n-7)}
 {n(3n-1)(2n-1)(3n-2)}\quad(n\geq1).
\]
This proves the recurrence recorded by Mathar and, since $T(0)=1$,
the stated generating function
\[
 \sum_{n\geq0}T(n)z^n=
 {}_4F_3(1/12,5/12,7/12,11/12;1/3,1/2,2/3;27648z).
\]
Finally Stirling's expansion gives Penson's asymptotic expression:
\[
 T(n)=\frac{27648^n}{\sqrt{12\pi n}}
 \left(1+\frac{1}{36n}+\frac{1}{2592n^2}+O(n^{-3})\right).
\]
Indeed the square-root factor is $1/\sqrt{12\pi n}$, and the
$1/n$ coefficient in the logarithm is
$(1/12+1-1/6-1/4-1/3)/12=1/36$; the logarithm has no $1/n^2$
term, so exponentiation gives $1/(2\cdot36^2)=1/2592$.
\end{proof}

The other A295431 questions have existing proofs.
Bajaj and Burns' August 2026 factorial-ratio theorem
\cite[Section~1]{tower}, followed by Blanco G\'omez's preprint of
17 October 2026 \cite[Corollary~1.4(d)]{blanco}, gives the full congruences
$T(np^k)\equiv T(np^{k-1})\pmod{p^{3k}}$ for $p\geq5$ and
$n,k\geq1$. Blanco G\'omez also proves the integrality of
$T(n/2)$ and $T(n/3)$ as Corollary~1.5 \cite{blanco}.
These two integrality questions already appear in Bala's 2016 note
\cite[p.~5]{bala2016}.
Its integer-argument congruences are attributed to the classical
Jacobsthal--Kazandzidis congruence; its Section~1 also identifies
Zudilin's earlier one-level result. These statements are distinct from
the product-divisibility families proved here.

\begingroup\small
\begin{longtable}{p{0.45\linewidth}p{0.46\linewidth}}
\toprule
Source item&Coverage\\
\midrule
A295462, References: binary-spacing conjecture&Proved in Theorem~\ref{thm:main}.\\
A295462, References: two divisibilities&Reproved; prior proof Bala, 2026 \cite{bala}; product and parity consequences in Corollary~\ref{cor:quotient}.\\
A295462, Formula: binomial, hypergeometric, asymptotic identities&Proved in Corollary~\ref{cor:formulas}, with the source attributions.\\
A048718, Name, Comments and Formula: sparse binary language and ordered recursion&Proved: the same index set and the counting/ordered recursions in Corollary~\ref{cor:count}.\\
A048718, Formula: all eleven Hanna XOR identities (2006)&Reproved; prior proof Karlsson, 2022 \cite{walnut}; written alternative proofs in Corollary~\ref{cor:xor}, including the masks 21 and 37.\\
A295431, Crossrefs: sporadic list, including row 32&Reproved; prior proof Bober, 2009 \cite[Table~2]{bober}, for the two displayed ratios' integrality; the classification is cited.\\
A295431, Comments (28 August 2025): five fixed divisibilities&Proved in Corollary~\ref{cor:families}; the 770 expression has exactly three denominator factors.\\
A295431, same Comments: $C(k,R)$ for $k=1,2,3$ and $D(R)$&Proved in Corollary~\ref{cor:families}, with explicit uniform constants.\\
A295431, Formula: two binomial identities, ${}_4F_3$, recurrence, asymptotic expansion&Proved in Corollary~\ref{cor:formulas}.\\
A295431, Formula: the explicit moment density and its shape/uniqueness&Not covered.\\
A295431, Comments (2020--2021): congruences modulo $p^{3k}$&Prior proofs cited: Bajaj--Burns (August 2026) \cite{tower}, then Blanco G\'omez (October 2026) \cite[Corollary~1.4(d)]{blanco}, via classical binomial scaling; Zudilin (2019) \cite[Section~5]{zudilin} covers $k=1$.\\
A295431, Comments (9 June 2026), and Bala (2016, p.~5): Gamma-factorial half/third arguments&Prior proof cited: Blanco G\'omez \cite[Corollary~1.5]{blanco}; earlier posing source Bala \cite{bala2016}.\\
A295431, Comments: analogous divisibility for the other 51 sporadic ratios&Not covered; their floor functions require separate analysis.\\
\bottomrule
\end{longtable}
\endgroup

\paragraph{Formal statement and finite checks.}
The declaration \texttt{result} in the pinned
\href{\leanurl}{Lean source file}
proves \texttt{claim} in the namespace
\nolinkurl{D5.S3.Arith.FactorialRatio.BalaA295462Parity}:
for every natural $n$ there is a natural $m$ equal, in the rationals,
to \eqref{eq:a}, and $m$ is odd if and only if
\texttt{spacedBits n} holds. That definition requires every pair of
occupied positions $i<j$ to satisfy $i+4\leq j$.
This is the integrality and parity assertion of Theorem~\ref{thm:main}.
Corollaries~\ref{cor:quotient}--\ref{cor:formulas}, including all eleven XOR
identities and the two A295431 product families, have written proofs.
The pinned declaration covers Theorem~\ref{thm:main} only.
The accompanying \texttt{verify.py} checks all 30 floor
intervals with rational arithmetic, the binary valuation criterion for
$0\leq n<2^{16}$, exact factorial values and product divisibility for
$0\leq n\leq256$, and the stated coefficient and counting identities.
It also checks the twelve cells for $G$, all eleven XOR identities,
the five A295431 divisibilities, and both product families for
$0\leq n\leq256$, $1\leq R\leq32$ and $k=1,2,3$.
Its finite ranges are separate from the all-input proofs.

\begin{thebibliography}{99}
\bibitem{oeis} The OEIS Foundation, \emph{A295462}, entry by Gheorghe
Coserea (28 November 2017), References contribution by Peter Bala
(23 August 2026), internal record \#19.
\url{https://oeis.org/A295462/internal}.
\bibitem{bala} Peter Bala, \emph{A295462(n) is divisible by 6*n+1 and
10*n+1}, 23 August 2026, 3 pp.
\url{https://oeis.org/A295462/a295462.pdf}.
\bibitem{bober} Jonathan W. Bober, \emph{Factorial ratios, hypergeometric
series, and a family of step functions}, J. London Math. Soc. (2)
\textbf{79} (2009), no.~2, 422--444.
\href{https://doi.org/10.1112/jlms/jdn078}{doi:10.1112/jlms/jdn078}.
Preprint \href{https://arxiv.org/abs/0709.1977v1}{arXiv:0709.1977v1},
Theorem~3.1, Lemma~3.2 and Table~2, row~32.
\bibitem{sound} K. Soundararajan, \emph{Integral factorial ratios},
arXiv:1901.05133v1, 16 January 2019, Theorem~1.1 and its proof.
\url{https://arxiv.org/abs/1901.05133v1}.
\bibitem{spaced} The OEIS Foundation, \emph{A048718}, entry by Antti
Karttunen (30 March 1999), Comments and Formula fields, internal record
\#25. \url{https://oeis.org/A048718/internal}.
\bibitem{walnut} Sebastian Karlsson, \emph{Walnut code that verifies the
conjectures of Paul D. Hanna}, linked from A048718, 31 December 2022.
\url{https://oeis.org/A048718/a048718.txt}.
\bibitem{sporadic} The OEIS Foundation, \emph{A295431}, entry by Gheorghe
Coserea (22 November 2017), Crossrefs and Comments fields, internal record
\#101. \url{https://oeis.org/A295431/internal}.
\bibitem{zudilin} Wadim Zudilin, \emph{Congruences for q-binomial
coefficients}, arXiv:1901.07843v2, 1 April 2019, Section~5, Theorem~3.
\url{https://arxiv.org/abs/1901.07843v2}.
\bibitem{tower} Ravi Bajaj and Alexander Burns, \emph{Balanced factorial
ratios inherit a cubic tower}, 2026, Section~1.
\href{https://github.com/rbajaj5/a183068-supercongruence/blob/d819eb8e8e70fdd386978adc4cbbdb2033466d40/related-results/BalancedFactorialRatioCubicTowers.md}{Pinned scientific text}.
\bibitem{bala2016} Peter Bala, \emph{Some integer ratios of factorials},
August 2016, 6 pp., p.~5.
\url{https://oeis.org/A276098/a276098.pdf}.
\bibitem{blanco} Roberto Blanco G\'omez, \emph{Supercongruences modulo
$p^{3k}$ and integrality for factorial ratios with fractional arguments,
with proofs of conjectures of P.~Bala}, Zenodo preprint, 17 October 2026,
8 pp., Corollaries~1.4(d) and 1.5, proofs in Section~5.2.
\href{https://doi.org/10.5281/zenodo.23426707}{doi:10.5281/zenodo.23426707}.
\end{thebibliography}
\noindent\small Copyright \textcopyright\ 2026 Haobo Ma and Wenlin Zhang.
Licensed under Creative Commons Attribution 4.0 International
(CC BY 4.0), \url{https://creativecommons.org/licenses/by/4.0/}.
\end{document}
